\documentclass[12pt]{scrartcl}
\usepackage[hagavi]{azzam}

\title{Pre Screening OSK SMA}
\author{120 menit}
\date{Sabtu, 10 Oktober 2026}

\begin{document}
\maketitle

\noindent Tes ini terdiri dari 16 soal isian singkat setingkat OSK SMA:
4 soal aljabar, 4 soal teori bilangan, 4 soal kombinatorika, dan 4 soal geometri.
Tuliskan jawaban akhirnya saja. Setiap jawaban yang benar bernilai 1 poin,
sedangkan jawaban yang salah atau kosong bernilai 0. Dilarang menggunakan kalkulator.

\section{Aljabar}
\begin{enumerate}
	% Diadaptasi dari OSK SMA 2023 no 5 (fungsi kuadrat)
	\item Diberikan fungsi kuadrat $f(x)=ax^{2}+bx+c$ yang memenuhi $f(3)=9$ dan $f(8)=64$.
	Jika $a \ne 1$, nilai dari $\dfrac{c-b}{a-1}$ adalah \dots

	% Diadaptasi dari OSK SMA 2024 no 6 (barisan geometri)
	\item Misalkan $u_{1},u_{2},u_{3},\dots$ suatu barisan geometri dengan $u_{1}>u_{2}$.
	Jika $u_{2}=9$ dan $u_{5}+u_{7}=\dfrac{10}{3}u_{6}$, nilai dari $u_{1}$ adalah \dots

	% Diadaptasi dari OSK SMA 2024 no 10 (manipulasi aljabar)
	\item Misalkan $x, y$ bilangan real positif dengan $x>y$.
	Diketahui $x^{2}+y^{2}=\dfrac{313}{156}xy$, maka $\dfrac{x+y}{x-y}$ adalah \dots

	% Diadaptasi dari OSK SMA 2025 no 9 (polinomial)
	\item Diketahui polinomial $P$ memenuhi $P(3^{b}+1)=3^{4b}+7$ untuk setiap bilangan asli $b$.
	Nilai dari $P(3)$ adalah \dots
\end{enumerate}

\section{Teori Bilangan}
\begin{enumerate}[resume]
	% Diadaptasi dari OSK SMA 2023 no 4 (bilangan kuadrat)
	\item Suatu bilangan empat digit $\overline{6ab9}$ merupakan suatu bilangan kuadrat.
	Nilai dari $a+b$ adalah \dots

	% Diadaptasi dari OSK SMA 2025 no 2 (faktorial)
	\item Bilangan bulat positif terkecil $n$ sehingga $n!$ habis dibagi $1540$ adalah \dots

	% Diadaptasi dari OSK SMA 2023 no 12 (persamaan Diophantine)
	\item Jika bilangan asli $x$ dan $y$ memenuhi
	\[ x(x-y)=7y-4 \]
	nilai dari $x+y$ adalah \dots

	% Diadaptasi dari OSK SMA 2023 no 8 (aritmetika modular)
	\item Sisa pembagian dari $5^{2025}+11^{2026}$ oleh $64$ adalah \dots
\end{enumerate}

\section{Kombinatorika}
\begin{enumerate}[resume]
	% Diadaptasi dari OSK SMA 2024 no 5 (kombinasi, jumlah digit)
	\item Bilangan KOTA adalah bilangan 4 angka yang tidak dimulai dengan angka 0 dan hasil
	penjumlahan semua digitnya sama dengan 9. Sebagai contoh, 2025 merupakan bilangan KOTA.
	Banyaknya bilangan KOTA adalah \dots

	% Diadaptasi dari OSK SMA 2025 no 5 (prinsip sarang merpati)
	\item Afif menuliskan sepuluh bilangan bulat positif yang lebih kecil dari 20.
	Ia memastikan bahwa penjumlahan dua bilangan mana pun di antara sepuluh bilangan tersebut
	tidak sama dengan 20. Bilangan positif yang pasti ditulis Afif adalah \dots

	% Diadaptasi dari OSK SMA 2023 no 2 (kombinasi)
	\item Di dalam suatu laci terdapat 8 pasang kaos kaki yang setiap pasangnya berbeda dengan
	pasangan lain. Diambil 6 kaos kaki sekaligus secara acak. Banyak cara pengambilan sehingga
	di antara yang terambil terdapat tepat sepasang kaos kaki yang berpasangan adalah \dots

	% Diadaptasi dari OSK SMA 2025 no 18 (pencacahan dasar)
	\item Suatu percobaan mengundi sebuah dadu beberapa kali dan percobaan berhenti setelah muncul
	mata dadu lebih kecil dari 5 sebanyak dua kali. Banyak kemungkinan percobaan berhenti pada
	pengundian ke-4 atau sebelumnya adalah \dots
\end{enumerate}

\section{Geometri}
\begin{enumerate}[resume]
	% Diadaptasi dari OSK SMA 2024 no 1 (segiempat)
	\item Sebuah persegi dibagi menjadi dua persegi panjang. Diketahui hasil penjumlahan kedua
	keliling persegi panjang tersebut adalah 66, maka luas persegi semula adalah \dots

	% Diadaptasi dari OSK SMA 2022 kode 2 (garis istimewa segitiga)
	\item Diberikan $ABC$ siku-siku sama kaki dengan panjang $BC=AB$ dan $L$ titik tengah $BC$.
	Titik $P$ pada sisi $AC$ sehingga $BP$ tegak lurus dengan $AL$.
	Jika panjang $CP=24\sqrt{2}$, panjang $AB$ adalah \dots

	% Diadaptasi dari OSK SMA 2022 kode 2 (lingkaran luar dan dalam)
	\item Diberikan segitiga siku-siku $ABC$ dengan luas 175. Misalkan $R$ adalah panjang jari-jari
	lingkaran luar segitiga $ABC$ dan $r$ adalah panjang jari-jari lingkaran dalam segitiga $ABC$.
	Jika $R+r=20$, panjang sisi miring dari segitiga $ABC$ adalah \dots

	% Diadaptasi dari OSK SMA 2026 no 17 (luas dan perbandingan)
	\item Diberikan segitiga $ABC$. Titik $D$ pada $AB$ sehingga $\dfrac{AD}{DB}=\dfrac{5}{2}$ dan
	titik $E$ pada $AC$ sehingga $\dfrac{AE}{EC}=\dfrac{1}{3}$. Misalkan $M$ adalah titik tengah $BC$
	dan $N$ adalah titik potong $AM$ dan $DE$. Jika luas dari segitiga $ABC$ adalah 1134,
	luas dari segitiga $ADN$ adalah \dots
\end{enumerate}

% ============================================================
% SOLUSI SINGKAT (mulai dari halaman baru).
% Untuk menyembunyikan solusi, beri tanda % di awal setiap baris
% dari blok ini sampai tepat sebelum \end{document}.
% ============================================================
% % \newpage
% % \section*{Solusi Singkat}
% \begin{enumerate}
% 	\item Kurangkan kedua persamaan: $55a+5b=55$, sehingga $b=11-11a$.
% 	Dari $f(3)=9$ diperoleh $c=9-9a-3b=24a-24$.
% 	Jadi $c-b=35a-35$ dan $\dfrac{c-b}{a-1}=\boxed{35}$.

% 	\item Tulis $u_{5}+u_{7}=u_{6}\left(\dfrac{1}{r}+r\right)$ dengan $r$ rasio barisan,
% 	sehingga $r+\dfrac{1}{r}=\dfrac{10}{3}$. Diperoleh $r=3$ atau $r=\dfrac{1}{3}$.
% 	Karena $u_{1}>u_{2}$, maka $r=\dfrac{1}{3}$ dan $u_{1}=\dfrac{9}{1/3}=\boxed{27}$.

% 	\item Misalkan $S=\dfrac{x^{2}+y^{2}}{xy}=\dfrac{313}{156}$. Maka
% 	\[ \left(\frac{x+y}{x-y}\right)^{2}=\frac{S+2}{S-2}
% 	=\frac{625/156}{1/156}=625. \]
% 	Karena $x>y>0$, nilai yang diminta adalah $\boxed{25}$.

% 	\item Misalkan $x=3^{b}+1$, sehingga $3^{b}=x-1$ dan $P(x)=(x-1)^{4}+7$.
% 	Jadi $P(3)=2^{4}+7=\boxed{23}$.

% 	\item Akar kuadratnya terletak di antara 78 dan 83. Agar kuadratnya berakhiran 9,
% 	akarnya harus berakhiran 3 atau 7, jadi akarnya 83 dan bilangannya $83^{2}=6889$.
% 	Nilai $a+b=8+8=\boxed{16}$.

% 	\item Karena $1540=2^{2}\cdot 5\cdot 7\cdot 11$, $n!$ harus memuat faktor prima 11,
% 	jadi $n \ge 11$. Nilai $11!$ memuat $2^{2}$, $5$, $7$, dan $11$, jadi $n=\boxed{11}$.

% 	\item Tulis $x^{2}+4=y(x+7)$, sehingga $y=x-7+\dfrac{53}{x+7}$.
% 	Jadi $x+7$ membagi 53. Karena 53 prima dan $x+7>7$, maka $x+7=53$.
% 	Diperoleh $x=46$, $y=40$, dan $x+y=\boxed{86}$.

% 	\item Modulo 64 berlaku $5^{16} \equiv 11^{16} \equiv 1$.
% 	Karena $2025 \equiv 9$ dan $2026 \equiv 10$ modulo 16, cukup hitung
% 	$5^{9}$ dan $11^{10}$. Dari $5^{4}=625 \equiv 49$ diperoleh
% 	$5^{8} \equiv 49^{2}=2401 \equiv 33$, sehingga $5^{9} \equiv 33\cdot 5=165 \equiv 37$.
% 	Dari $11^{2}=121 \equiv 57$ diperoleh $11^{4} \equiv 57^{2}=3249 \equiv 49$ dan
% 	$11^{8} \equiv 33$, sehingga $11^{10} \equiv 33\cdot 57=1881 \equiv 25$.
% 	Jumlahnya $37+25=\boxed{62}$.

% 	\item Tulis $a+b+c+d=9$ dengan $a \ge 1$. Misalkan $a'=a-1$, maka
% 	$a'+b+c+d=8$ dengan semua peubah $\ge 0$. Banyak solusinya
% 	$\binom{8+4-1}{4-1}=\binom{11}{3}=\boxed{165}$.

% 	\item Pasangan bilangan yang berjumlah 20 adalah $(1,19)$, $(2,18)$, \dots, $(9,11)$,
% 	serta $10+10$. Sepuluh bilangan yang dipilih tidak boleh memuat dua bilangan yang
% 	berpasangan, jadi tepat satu bilangan diambil dari tiap pasangan dan bilangan 10
% 	wajib ikut. Bilangan yang pasti ditulis adalah $\boxed{10}$.

% 	\item Pilih pasangan yang utuh: 8 cara. Dari 7 pasang sisanya, pilih 4 pasang:
% 	$\binom{7}{4}=35$ cara, lalu dari tiap pasang pilih sebelah kiri atau kanan: $2^{4}=16$ cara.
% 	Totalnya $8\cdot 35\cdot 16=\boxed{4480}$.

% 	\item Sebut $S$ kejadian muncul mata $<5$ (4 cara) dan $G$ sisanya (2 cara).
% 	Percobaan yang berhenti pada undian ke-$n$ berbentuk $n-1$ undian berisi tepat satu $S$
% 	diakhiri $S$: ada $(n-1)\cdot 4^{2}\cdot 2^{n-2}$ kemungkinan.
% 	Untuk $n=2,3,4$ diperoleh $16+64+192=\boxed{272}$.

% 	\item Misalkan sisi persegi $s$ dan potongannya sejajar salah satu sisi.
% 	Jumlah kedua keliling adalah $2(s+h_{1})+2(s+h_{2})=4s+2s=6s$ karena $h_{1}+h_{2}=s$.
% 	Dari $6s=66$ diperoleh $s=11$, jadi luasnya $\boxed{121}$.

% 	\item Letakkan $B=(0,0)$, $C=(s,0)$, $A=(0,s)$ dengan $s=AB$.
% 	Maka $L=(s/2,0)$, garis $BP$ adalah $y=x/2$, dan garis $AC$ adalah $x+y=s$,
% 	sehingga $P=(2s/3,s/3)$. Jadi $CP=s\sqrt{2}/3=24\sqrt{2}$, maka $s=\boxed{72}$.

% 	\item Misalkan sisi miring $c=2R$ dan luas $L=ab/2=175$.
% 	Dari $r=\dfrac{a+b-c}{2}$ diperoleh $(2r+c)^{2}=c^{2}+4L$, yaitu $r^{2}+rc=L$.
% 	Dengan $c=40-2r$, maka $r(40-r)=175$, yaitu $r^{2}-40r+175=0$.
% 	Solusinya $r=5$ atau $r=35$ (tidak mungkin), jadi $c=40-10=\boxed{30}$.

% 	\item Tulis $D=\dfrac{2A+5B}{7}$ dan $E=\dfrac{3A+C}{4}$.
% 	Titik $N=(1-t)D+tE$ pada $AM$ berbentuk $\lambda A+(1-\lambda)\dfrac{B+C}{2}$.
% 	Menyamakan koefisien $B$ dan $C$: $\dfrac{5(1-t)}{7}=\dfrac{t}{4}$,
% 	sehingga $t=\dfrac{20}{27}$. Koefisien $C$ dari $N$ adalah $\dfrac{5}{27}$, jadi
% 	\[ \frac{[ADN]}{[ABC]}=\frac{AD}{AB}\cdot\frac{5}{27}
% 	=\frac{5}{7}\cdot\frac{5}{27}=\frac{25}{189}. \]
% 	Maka $[ADN]=1134\cdot\dfrac{25}{189}=\boxed{150}$.
% \end{enumerate}

\end{document}
